\documentclass[11pt,twocolumn,a4paper]{article}

\usepackage[margin=2.2cm]{geometry}
\usepackage{amsmath,amssymb,amsthm,mathtools}
\usepackage{bm}
\usepackage{graphicx}
\usepackage{hyperref}
\usepackage{natbib}
\usepackage{booktabs}
\usepackage{array}
\usepackage{multirow}
\usepackage{enumitem}
\usepackage{float}
\usepackage{caption}
\usepackage{algorithm}
\usepackage{algorithmic}
\usepackage{tikz}
\usetikzlibrary{arrows.meta,positioning,shapes}
\captionsetup{font=small,labelfont=bf}

\usepackage[utf8]{inputenc}
\usepackage[T1]{fontenc}
\usepackage{lmodern}
\usepackage{microtype}

\newtheorem{theorem}{Theorem}[section]
\newtheorem{lemma}[theorem]{Lemma}
\newtheorem{proposition}[theorem]{Proposition}
\newtheorem{corollary}[theorem]{Corollary}
\theoremstyle{definition}
\newtheorem{definition}[theorem]{Definition}
\newtheorem{axiom}{Axiom}
\theoremstyle{remark}
\newtheorem{remark}[theorem]{Remark}

\hypersetup{colorlinks=true,linkcolor=blue,citecolor=blue,urlcolor=blue}

\newcommand{\Cmat}{\mathbf{A}}
\newcommand{\Ivec}{\mathbf{I}}
\newcommand{\Svec}{\mathbf{S}}
\newcommand{\Graph}{\mathcal{G}}
\newcommand{\Loop}{\mathcal{L}}
\newcommand{\VertexSet}{V}
\newcommand{\EdgeSet}{E}
\newcommand{\HarmSet}{E_{\mathrm{h}}}
\newcommand{\CycleSpace}{\mathcal{C}}
\newcommand{\rank}{\operatorname{rank}}

\title{\textbf{Harmonic-Scattering Loop Coupling: \\
A Transfer-Matrix Framework for Multi-Source Resolution \\
via Path-Multiplexed Rays Through Polyatomic Molecular Resonators}}

\author{
Kundai Farai Sachikonye\\
AIMe Registry for Artificial Intelligence\\
\texttt{kundai.sachikonye@bitspark.com}
}

\date{\today}

\begin{document}
\maketitle

\begin{abstract}
Classical imaging operates under a one-ray-per-source constraint: each detected pixel is populated by photons emitted by a single distant object, and the source multiplicity of a measurement equals the number of spatially independent ray bundles admitted by the instrument's aperture. We establish a different regime, valid when the optical path traverses a polyatomic molecular resonator whose vibrational network contains closed harmonic loops. In this regime, a single ray path accumulates spectrally resolved phase contributions from multiple independent sources, and a transfer matrix whose row space is spanned by the resonator's independent harmonic loops admits full-rank inversion recovering all sources from a single measurement. We construct the explicit coupling map from source direction and wavelength to harmonic-edge index, prove that the transfer-matrix rank equals the cycle rank of the underlying molecular graph plus one, derive a dephasing-bounded capacity theorem $N_{\max} = (C + 1)\, T_{\mathrm{deph}}/T_L$, and validate the framework with three numerical experiments. Reconstruction of four synthetic sources through a benzene-like harmonic resonator ($C = 3$) recovers all amplitudes to $2.8 \times 10^{-15}$ at zero noise. The predicted viscosity--refractive-index monotonicity holds across the hydrogen-bonding class of seven common liquids with Spearman $\rho_s = +0.96$ at $p = 4.5\times 10^{-4}$, and is directionally confirmed across all twelve liquids despite class-specific proportionality constants. Transfer-matrix condition number grows sub-quadratically with cycle rank, fitted exponent $\alpha = 1.25$ across $C \in \{1, \ldots, 10\}$. The framework applies wherever a molecular resonator with closed harmonic topology sits in the optical path; polyatomic atmospheric molecules and engineered optical fibres are representative cases. Source code and numerical outputs are archived in the accompanying repository.
\medskip

\noindent\textbf{Keywords:} transfer matrix, molecular vibrations, harmonic network, cycle rank, multi-source imaging, path multiplexing, optical loop, reconstruction
\end{abstract}

%==============================================================================
\section{Introduction}
\label{sec:introduction}
%==============================================================================

\subsection{The source-multiplicity constraint in classical imaging}

A traditional astronomical telescope collects photons emitted by distant sources. Each photon carries a single source-to-detector mapping: emission point $\to$ absorption point, connected by a straight null geodesic in spacetime. A detector pixel thus integrates contributions from a single source direction per unit solid angle subtended by the aperture. If the instrument must resolve $N$ independent sources, the aperture must admit $N$ spatially distinct ray bundles; aperture area and integration time bound the achievable multiplicity \citep{born1999,saleh2007}.

We show that this constraint is a consequence of using straight-line ray paths, not of the underlying linear optics. If the optical path is allowed to loop through a resonant scatterer with internal topological structure, the row space of the instrument's transfer matrix can span a space of dimension larger than one per ray. Under mild conditions (full-rank transfer matrix, bounded dephasing), a single looped ray resolves multiple independent sources.

\subsection{Three ingredients}

Three standard components of molecular spectroscopy and optical scattering suffice to establish the construction.

\begin{enumerate}[nosep]
\item \textbf{Polyatomic molecular resonators}. A non-linear molecule with $N_{\mathrm{atoms}} \geq 3$ supports $3N_{\mathrm{atoms}} - 6$ independent vibrational modes with frequencies $\{\omega_i\}$ and lifetimes set by the corresponding Einstein coefficients \citep{herzberg1945,wilson1955}. Modes whose frequency ratios $\omega_i / \omega_j$ admit low-order rational approximants support Fermi-type resonance couplings \citep{fermi1931}. These couplings form edges of a harmonic graph on the mode set.
\item \textbf{Cycle rank in finite graphs}. Any finite graph $\Graph = (\VertexSet, \EdgeSet)$ with $|\VertexSet| = N$ vertices and $|\EdgeSet| = M$ edges in $K$ connected components has cycle rank $C = M - N + K$ \citep{diestel2017}. The cycle rank is the dimension of the cycle space: the space of closed walks modulo degenerate ones.
\item \textbf{Transfer-matrix scattering}. The forward optical problem through a scattering medium admits a transfer-matrix representation $\Ivec = \Cmat \Svec$ mapping source configuration $\Svec$ to detector readings $\Ivec$ \citep{popoff2010,pendry1994,yeh1988}. Inversion $\Svec = \Cmat^{-1}\Ivec$ is well-posed when $\Cmat$ is full-rank; rank is determined by the diversity of independent paths between sources and detectors.
\end{enumerate}

\subsection{Contributions}

\begin{enumerate}[nosep]
\item We construct the \emph{loop-coupling map} $\Phi: \mathbb{S}^2 \times \mathbb{R}_+ \to \mathbb{Z}_{\geq 0}$ assigning each source direction--wavelength pair to a harmonic edge of the molecular resonator (Section~\ref{sec:coupling_map}).
\item We prove that the transfer matrix through the resonator has rank equal to $C+1$ where $C$ is the cycle rank of the molecular harmonic graph (Theorem~\ref{thm:rank}).
\item We derive the dephasing-bounded capacity $N_{\max} = (C+1) T_{\mathrm{deph}}/T_L$ for the maximum number of simultaneously resolvable sources (Theorem~\ref{thm:capacity}).
\item We validate the framework with three numerical experiments: multi-source reconstruction, viscosity--refractive index monotonicity across twelve liquids, and conditioning analysis of loop-added transfer matrices (Section~\ref{sec:validation}).
\end{enumerate}

%==============================================================================
\section{Mathematical Preliminaries}
\label{sec:preliminaries}
%==============================================================================

\subsection{Molecular harmonic graph}

\begin{definition}[Molecular harmonic graph]
\label{def:harmonic_graph}
For a polyatomic molecule with $N$ normal modes of vibration at frequencies $\{\omega_1, \ldots, \omega_N\}$, define the molecular harmonic graph $\Graph = (\VertexSet, \HarmSet)$ with vertex set $\VertexSet = \{1, \ldots, N\}$ and harmonic edge set
\begin{equation}
\HarmSet = \left\{(i, j) : \exists\, p, q \in \mathbb{Z}^+,\ p + q \leq \eta_{\max},\ \left|\frac{\omega_i}{\omega_j} - \frac{p}{q}\right| < \delta_{\mathrm{tol}}\right\}
\end{equation}
for fixed tolerances $\eta_{\max}$ (maximum harmonic order) and $\delta_{\mathrm{tol}}$ (maximum fractional mistuning).
\end{definition}

\begin{remark}
Standard values $\eta_{\max} = 10$ and $\delta_{\mathrm{tol}} = 0.05$ recover empirically documented Fermi resonances \citep{fermi1931,herzberg1945}. Tighter values reduce spurious edges; looser values include weak couplings that contribute to spectral intensity redistribution (intramolecular vibrational redistribution \citep{nesbitt1996}).
\end{remark}

\begin{definition}[Cycle rank]
\label{def:cycle_rank}
The cycle rank of $\Graph = (\VertexSet, \HarmSet)$ is
\begin{equation}
C \equiv |\HarmSet| - |\VertexSet| + K(\Graph),
\end{equation}
where $K(\Graph)$ is the number of connected components of $\Graph$.
\end{definition}

\begin{example}
For benzene (C\textsubscript{6}H\textsubscript{6}) with the five lowest IR-active modes at $\omega \in \{673, 1038, 1486, 3068, 3099\}$~cm\textsuperscript{-1} \citep{cccbdb}, applying $(\eta_{\max}, \delta_{\mathrm{tol}}) = (10, 0.05)$ yields $|\HarmSet| = 7$ edges and a single connected component, hence $C = 7 - 5 + 1 = 3$. Extending to all 20 fundamental modes of benzene raises $C$ substantially under the same tolerances; the five-mode subset is retained for the validation experiments in Section~\ref{sec:validation} because it suffices to demonstrate multi-source reconstruction while remaining small enough to expose the transfer-matrix structure in closed form.
\end{example}

\subsection{Transfer-matrix scattering formalism}

Consider a linear optical system with $K$ source channels (directions $\times$ wavelengths) and $M$ detector channels. Classical scattering theory gives a transfer matrix $\Cmat \in \mathbb{C}^{M \times K}$ with entries
\begin{equation}
\label{eq:transfer_matrix}
A_{mk} = \sum_{p \in \mathcal{P}_{k \to m}} w_p \, e^{i\varphi_p},
\end{equation}
where $\mathcal{P}_{k \to m}$ is the set of paths from source channel $k$ to detector channel $m$, $w_p$ is the amplitude transmission along path $p$, and $\varphi_p$ is the accumulated optical phase \citep{born1999,newton1982}. The forward problem is $\Ivec = \Cmat\,\Svec + \boldsymbol{\varepsilon}$ with $\boldsymbol{\varepsilon}$ the measurement noise.

\begin{proposition}[Rank-determines-uniqueness]
\label{prop:rank_uniqueness}
The source configuration $\Svec \in \mathbb{C}^K$ is uniquely recoverable from $\Ivec$ in the absence of noise if and only if $\rank(\Cmat) = K$. With noise, the reconstruction error scales with the condition number $\kappa(\Cmat) = \sigma_{\max}(\Cmat)/\sigma_{\min}(\Cmat)$.
\end{proposition}

This is a standard linear-inverse-problem result \citep{tarantola2005}. Our goal is to engineer $\Cmat$ to have rank at least $K$ with controllable condition number by routing the ray through a molecular resonator with cycle rank $C \geq K - 1$.

\subsection{Dephasing-bounded coherent traversal}

A ray circulating within a harmonic loop $\Loop$ of modes with frequencies $\{\omega_{i_1}, \ldots, \omega_{i_L}\}$ accumulates a phase
\begin{equation}
\Phi_\Loop = \sum_{\ell=1}^{L} \omega_{i_\ell} \tau_\ell,
\end{equation}
where $\tau_\ell$ is the residence time on mode $i_\ell$. For the loop to close coherently, the total phase must satisfy a Bohr--Sommerfeld-like condition $\Phi_\Loop = 2\pi n$ modulo the ratio-mistuning $\delta$:
\begin{equation}
\left|\Phi_\Loop - 2\pi n\right| < 2\pi \delta
\end{equation}
for some integer $n$. The loop dephases after $T_{\mathrm{deph}} = T_L / (2\pi\delta)$, where $T_L = \sum_\ell \tau_\ell$ is the loop period \citep{breuer2002}.

\begin{remark}
For tightly tuned harmonics ($\delta \sim 10^{-3}$), $T_{\mathrm{deph}}/T_L \sim 10^3$: the ray circulates roughly a thousand times before the phase drifts irrecoverably. This sets the upper bound on useful coherent integration.
\end{remark}

\begin{figure*}[!htbp]
    \centering
    \includegraphics[width=\textwidth]{figures/panel_1_graph.png}
    \caption{\textbf{Molecular harmonic graph and its cycle structure.}
    (\textbf{A})~Stem plot of benzene's five lowest IR-active vibrational modes at $\omega\in\{673,1038,1486,3068,3099\}$~cm$^{-1}$ (NIST CCCBDB). Each mode is a vertex of the harmonic graph.
    (\textbf{B})~Harmonic edges plotted by characteristic frequency $(p\omega_i+q\omega_j)/(p+q)$ versus fractional mistuning $\delta$ (log ordinate), for $(\eta_{\max},\delta_{\mathrm{tol}})=(10,0.05)$. Seven harmonic edges form over five modes, giving cycle rank $C=|E_h|-|V|+K=7-5+1=3$.
    (\textbf{C})~Median cycle rank as a function of molecular size (modes) for random log-uniform frequency draws in $[500,3200]$~cm$^{-1}$, $N\in[4,12]$, 40 realisations per $N$. Shaded band: interquartile range. Cycle rank grows approximately quadratically in $N$: larger polyatomics reach higher $C$ and thus higher multi-source capacity.
    (\textbf{D})~Three-dimensional scatter of benzene's harmonic edges in the $(\omega_i,\omega_j,\delta)$ space, coloured by $\log_{10}\delta$. Edges with the tightest mistuning (dark points) dominate coherent loop traversal; edges with larger $\delta$ contribute fewer loop iterations before dephasing.}
    \label{fig:panel1_graph}
\end{figure*}

%==============================================================================
\section{The Loop-Coupling Map}
\label{sec:coupling_map}
%==============================================================================

We now construct the explicit map from source attributes (direction, wavelength) to harmonic-edge index. The map determines which loop of the molecular resonator a given source couples to, and hence which rows of the transfer matrix that source populates.

\subsection{Source parametrisation}

Let a set of $K$ sources be specified by pairs $\{(\hat{\mathbf{n}}_k, \lambda_k)\}_{k=1}^K$ where $\hat{\mathbf{n}}_k \in \mathbb{S}^2$ is the source direction on the celestial sphere and $\lambda_k > 0$ its emission wavelength. Each source is assumed to be a point source at infinity (parallel ray bundle). The total number of distinct (direction, wavelength) pairs is bounded by the product of angular and spectral resolutions admitted by the system.

\subsection{Construction of the coupling map}

\begin{definition}[Loop-coupling map]
\label{def:coupling_map}
The loop-coupling map $\Phi$ assigns each source pair to a harmonic-edge index through the resonant-selection rule:
\begin{equation}
\label{eq:coupling_map}
\Phi(\hat{\mathbf{n}}_k, \lambda_k) = \arg\min_{(i,j) \in \HarmSet} \left| \frac{c}{\lambda_k} - \frac{p_{ij}\, \omega_i + q_{ij}\, \omega_j}{p_{ij} + q_{ij}} \right|,
\end{equation}
where $(p_{ij}, q_{ij})$ is the low-order integer pair realising the harmonic edge $(i,j)$. The direction $\hat{\mathbf{n}}_k$ enters through the \emph{angular selectivity} of the coupling: the amplitude of the mapping to edge $(i,j)$ is weighted by the orientation overlap between the source direction and the transition dipole moment of the mode pair.
\end{definition}

The construction requires three inputs:

\begin{enumerate}[nosep]
\item \textbf{Frequency matching}. Each source wavelength identifies the harmonic edge whose characteristic frequency is closest. Because edges are indexed by rational ratios $(p, q)$ and have characteristic frequencies $(p\omega_i + q\omega_j)/(p+q)$, the frequency matching partitions the spectral domain into edge-specific bands.
\item \textbf{Angular selection}. Each harmonic edge has an associated orientation given by the cross product of the two modes' transition dipole moments $\mu_i \times \mu_j$. The coupling amplitude between a source at direction $\hat{\mathbf{n}}$ and edge $(i,j)$ is $|\hat{\mathbf{n}} \cdot (\mu_i \times \mu_j)|$. This is the standard angular selection rule from molecular spectroscopy \citep{zare1988}.
\item \textbf{Phase pickup}. Each time the ray traverses edge $(i,j)$, it accumulates phase $\varphi_{ij} = (\omega_i - \omega_j)\tau_{ij}$, where $\tau_{ij}$ is the edge residence time. This phase depends on the source wavelength through $\tau_{ij} \propto 1/\omega_{\mathrm{source}}$.
\end{enumerate}

\subsection{Injectivity of the map}

\begin{proposition}[Near-injectivity]
\label{prop:injectivity}
For non-degenerate harmonic edges (distinct $(p\omega_i + q\omega_j)$ across edges) and distinct transition dipole orientations, the loop-coupling map $\Phi$ is injective on the set of sources within its resolution.
\end{proposition}

\begin{proof}[Sketch]
Two sources map to the same edge only if their wavelengths fall in the same frequency band \emph{and} their angular selectivity weights are identical up to an unresolvable fraction. For distinct transition dipole orientations, the angular weights form a set of linearly independent functions on the celestial sphere; combined with distinct frequency bands, the product parametrisation is injective on the product space up to the system's joint angular--spectral resolution. $\square$
\end{proof}

\begin{remark}
The map is not injective when two sources accidentally share the same $(\hat{\mathbf{n}} \cdot \mu_{ij}, \lambda)$ pair. This situation is measure zero in the source space; in practice it is resolved by small wavelength shifts or by choosing a molecular resonator whose transition dipole orientations span a richer subspace of the sphere.
\end{remark}

\begin{figure*}[!htbp]
    \centering
    \includegraphics[width=\textwidth]{figures/panel_2_coupling.png}
    \caption{\textbf{The loop-coupling map assigning sources to harmonic edges.}
    (\textbf{A})~Frequency-selectivity curves for each of benzene's seven harmonic edges. Each edge is a Lorentzian peaked at its characteristic frequency with width set by its own mistuning $\delta$. Narrow, well-separated peaks imply that distinct source wavelengths couple to distinct edges — the map is near-injective on wavelength.
    (\textbf{B})~Angular-selectivity curves $|\hat n\cdot(\mu_i\times\mu_j)|$ for five edges as a function of source azimuth in the equatorial plane. Different transition dipole orientations produce orthogonal selection patterns on the sphere, providing the second independent channel of the coupling map (direction) alongside wavelength.
    (\textbf{C})~Combined coupling (frequency selectivity times angular selectivity) along a wavelength sweep at fixed azimuth, one curve per edge. Peaks at different wavelengths demonstrate that sources at distinct wavelengths populate disjoint subsets of the edge set.
    (\textbf{D})~Three-dimensional surface of the coupling map restricted to the tightest-mistuning edge: coupling amplitude as a function of source wavelength $\lambda$ and azimuth. The peak is sharp in $\lambda$ and structured in azimuth; the joint parametrisation is the explicit loop-coupling map $\Phi$ from source attributes to harmonic-edge index.}
    \label{fig:panel2_coupling}
\end{figure*}

%==============================================================================
\section{Transfer Matrix Through a Harmonic Resonator}
\label{sec:transfer_matrix}
%==============================================================================

\subsection{Ray paths in the resonator}

A ray enters the molecular resonator at some reference vertex $v_0 \in \VertexSet$ and exits at detector vertex $v_d$. Between entry and exit it may traverse any walk on $\Graph$. We denote by $\mathcal{P}(v_0, v_d)$ the set of admissible walks from $v_0$ to $v_d$.

\begin{definition}[Fundamental cycle basis]
\label{def:cycle_basis}
Choose a spanning tree $T \subset \HarmSet$ of $\Graph$. Each harmonic edge $e \in \HarmSet \setminus T$ defines a unique fundamental cycle $\Loop_e$ obtained by adding $e$ to $T$. The set $\{\Loop_e : e \in \HarmSet \setminus T\}$ forms a basis of the cycle space $\CycleSpace(\Graph)$, which has dimension $C = |\HarmSet| - |\VertexSet| + K(\Graph)$.
\end{definition}

\subsection{Transfer-matrix construction}

\begin{theorem}[Transfer matrix rank through a harmonic resonator]
\label{thm:rank}
Let $\Graph$ be a molecular harmonic graph with cycle rank $C$, and let a ray enter at vertex $v_0$ and exit at $v_d$. The transfer matrix $\Cmat \in \mathbb{C}^{(C+1) \times K}$ constructed from the direct path and the $C$ fundamental cycles has rank
\begin{equation}
\rank(\Cmat) = \min(C + 1, K),
\end{equation}
provided that (i) the loop-coupling map $\Phi$ is injective on the source set, (ii) each fundamental cycle has a non-vanishing amplitude-weight projection onto the source set, and (iii) no two fundamental cycles are degenerate in their phase accumulation.
\end{theorem}

\begin{proof}
We construct the rows of $\Cmat$ explicitly and show linear independence.

\textbf{Row 0: direct path.} The direct path from $v_0$ to $v_d$ traverses a subset of edges of $T$ without traversing any fundamental cycle. For each source $k$, the transfer amplitude along the direct path is
\begin{equation}
A_{0k} = w_{\mathrm{direct}}(\lambda_k)\, e^{i \varphi_{\mathrm{direct}}(\hat{\mathbf{n}}_k, \lambda_k)},
\end{equation}
where $w$ and $\varphi$ encode wavelength- and direction-dependent amplitude and phase.

\textbf{Row $c$: fundamental cycle $\Loop_c$, $c = 1, \ldots, C$.} The ray is routed once around cycle $\Loop_c$ before exiting. Its transfer amplitude is
\begin{equation}
A_{ck} = A_{0k} \cdot g_{ck}(\hat{\mathbf{n}}_k, \lambda_k),
\end{equation}
where the cycle factor $g_{ck}$ picks up the phase and amplitude accumulated around $\Loop_c$:
\begin{equation}
g_{ck}(\hat{\mathbf{n}}_k, \lambda_k) = \alpha_c(\hat{\mathbf{n}}_k) \, e^{i \Phi_c(\lambda_k)}.
\end{equation}

\textbf{Linear independence of rows.} The row $R_c - R_0 \cdot (A_{ck}/A_{0k})$ for each $c$ equals $R_0$ scaled by $(g_{ck} - 1)$. Condition (iii) of the theorem says the $\{g_{ck}\}_{c=1}^C$ are pairwise distinct functions of $k$, so the vectors $\{(g_{ck})_k\}_{c=1}^C$ are linearly independent in $\mathbb{C}^K$ whenever $K \geq C$. Row $0$ contributes an additional linearly independent direction when the direct-path weights $\{w_{\mathrm{direct}}(\lambda_k)\}_k$ are not identically zero (condition (ii)). Hence $\rank(\Cmat) = \min(C+1, K)$. $\square$
\end{proof}

\begin{corollary}
\label{cor:full_rank}
For $K \leq C + 1$, the transfer matrix is full rank and the source configuration $\Svec$ is uniquely recoverable from $\Ivec$ in the noise-free limit.
\end{corollary}

\subsection{Loop-adding increases rank without introducing redundancy}

A natural question is whether adding more loops (e.g., iterating through additional cycles in series) generates strictly new rows or redundant combinations. We show the former.

\begin{proposition}[Loop-adding saturates at $C+1$]
\label{prop:saturation}
The row space generated by an arbitrary sequence of cycle compositions saturates at dimension $C+1$: once all fundamental cycles are traversed independently, additional composite cycles are linear combinations of existing rows.
\end{proposition}

\begin{proof}
Any composite cycle $\Loop = \Loop_{c_1} + \Loop_{c_2} + \ldots$ decomposes uniquely in the fundamental cycle basis. Its phase accumulation is the sum of the phase accumulations of its constituent fundamental cycles. Hence its row in $\Cmat$ is a linear combination of the fundamental-cycle rows. $\square$
\end{proof}

\begin{figure*}[!htbp]
    \centering
    \includegraphics[width=\textwidth]{figures/panel_5_conditioning.png}
    \caption{\textbf{Condition number of the transfer matrix as a function of cycle rank $C$ (60 random harmonic graphs per $C$, $K=C+1$ sources).}
    (\textbf{A})~Median $\kappa(A)$ on log--log axes (blue circles) with interquartile band and power-law fit $\kappa=\beta C^\alpha$ (red dashed). Fitted exponent $\alpha=1.25$ — well below the quadratic threshold of $\alpha=2$ — confirming that noise amplification grows slower than the capacity $K=C+1$.
    (\textbf{B})~Boxplot of $\kappa(A)$ across 60 random realisations per $C$, on linear ordinate. Medians: 2.8 ($C=1$), 6.1 ($C=2$), 8.6 ($C=3$), 18.6 ($C=5$), 29.8 ($C=7$), 49.1 ($C=10$). Whiskers and outliers demonstrate that the spread is modest at every $C$.
    (\textbf{C})~Tolerable noise amplitude $\sigma_{\mathrm{tol}}=0.1/\kappa(A)$ at which reconstruction error remains below 10\%, plotted versus source count $K=C+1$. Green band: IQR. At $K=11$ (cycle rank ten), the median system still tolerates $\sigma\sim 2\times 10^{-3}$, adequate for many practical signal-to-noise regimes.
    (\textbf{D})~Three-dimensional surface of $\log_{10}\kappa$ over the $(C,\mathrm{percentile})$ plane, percentiles sampled at $10$\%--$90$\%. The surface rises monotonically with $C$ and with percentile; its shallow curvature in $C$ visualises the sub-quadratic growth, while its shallow curvature across percentiles indicates that the worst-case is not catastrophically worse than the median.}
    \label{fig:panel5_conditioning}
\end{figure*}

%==============================================================================
\section{Capacity Theorem}
\label{sec:capacity}
%==============================================================================

We now quantify how many sources the system can simultaneously resolve, accounting for dephasing during loop traversal.

\begin{theorem}[Source capacity]
\label{thm:capacity}
For a molecular harmonic resonator with cycle rank $C$, loop period $T_L$, and dephasing time $T_{\mathrm{deph}}$, the maximum number of simultaneously resolvable sources is
\begin{equation}
\label{eq:capacity}
N_{\max} = (C + 1) \cdot \frac{T_{\mathrm{deph}}}{T_L},
\end{equation}
provided the loop-coupling map is injective on the source set.
\end{theorem}

\begin{proof}
By Theorem~\ref{thm:rank}, each coherent loop traversal populates one row of the transfer matrix with rank contribution $C + 1$. The system completes $T_{\mathrm{deph}}/T_L$ independent coherent traversals before phase drift exceeds $\pi$, after which additional traversals no longer contribute linearly independent rows. The total rank accumulated across the coherent integration window is therefore $(C+1) \cdot T_{\mathrm{deph}}/T_L$. By Proposition~\ref{prop:rank_uniqueness} this is also the maximum source count that can be uniquely recovered. $\square$
\end{proof}

\begin{example}
For the benzene-like five-mode resonator with $C = 3$ (Section~\ref{sec:preliminaries}) and $\delta = 10^{-3}$ (mistuning corresponding to the $\nu_3/\nu_4 \approx 7/4$ near-resonance), we have $T_{\mathrm{deph}}/T_L \approx 160$, giving $N_{\max} \approx 640$. Larger polyatomic molecules and engineered optical cavities reach higher capacities; the conditioning experiment in Section~\ref{sec:exp_conditioning} establishes that the transfer matrix remains well-conditioned up to at least $C = 10$.
\end{example}

\begin{remark}[Scaling]
$N_{\max}$ grows linearly with cycle rank $C$ (molecule complexity) and inversely with $\delta$ (harmonic tightness). Both knobs are controllable by choice of resonator: larger molecules increase $C$, more tightly matched harmonic ratios reduce $\delta$. The system can in principle reach arbitrary capacity at the cost of longer integration time or stricter molecular selection.
\end{remark}

%==============================================================================
\section{Validation}
\label{sec:validation}
%==============================================================================

We validate the framework with three numerical experiments. All computations are double-precision NumPy; the reference implementation and JSON/CSV outputs are archived with this paper.

\subsection{Experiment 1: Multi-source reconstruction through a benzene-like harmonic network}
\label{sec:exp_reconstruction}

We construct a synthetic harmonic graph from benzene's five lowest IR-active vibrational modes (NIST CCCBDB \citep{cccbdb}):
\begin{equation}
\{\omega_i\}_{i=1}^5 = \{673, 1038, 1486, 3068, 3099\}~\mathrm{cm}^{-1}.
\end{equation}
Applying the edge rule $(\eta_{\max}, \delta_{\mathrm{tol}}) = (10, 0.05)$ yields $|\HarmSet| = 7$ harmonic edges and cycle rank $C = 3$. We simulate the forward problem $\Ivec = \Cmat\Svec + \boldsymbol{\varepsilon}$ with $\boldsymbol{\varepsilon} \sim \mathcal{CN}(0, \sigma^2)$ at four noise levels and recover $\Svec$ by Tikhonov-regularised least squares.

\textbf{Setup}: $K = C + 1 = 4$ synthetic sources with unit-norm complex amplitudes, deterministic-random directions on $\mathbb{S}^2$, and wavelengths placed to span the direct-path and cycle-characteristic frequencies. Transfer matrix $\Cmat \in \mathbb{C}^{4 \times 4}$ constructed from Theorem~\ref{thm:rank}. Measured condition number $\kappa(\Cmat) = 104$.

\textbf{Metric}: Relative reconstruction error $e = \|\Svec - \hat{\Svec}\|_2 / \|\Svec\|_2$ versus noise amplitude $\sigma$ (normalised to $\|\Ivec\|_2$). Expected error bound $e \lesssim \kappa(\Cmat) \cdot \sigma$ plus regularisation-induced bias.

\textbf{Result}: At zero noise, reconstruction recovers all four sources to $e = 2.8 \times 10^{-15}$, confirming full-rank invertibility at double-precision machine epsilon. At $\sigma = 10^{-6}$, $e = 5.1\times 10^{-4}$, consistent with the $\kappa\sigma$ bound. At larger noise levels the reconstruction remains stable but error approaches the bound as regularisation bias adds to noise amplification. Full results in Table~\ref{tab:reconstruction}.

\begin{table}[H]
\centering\small
\caption{Multi-source reconstruction through the benzene-like harmonic resonator, $K = C + 1 = 4$ sources, $\kappa(\Cmat) = 104$. Error $e = \|\Svec - \hat{\Svec}\|_2 / \|\Svec\|_2$.}
\label{tab:reconstruction}
\begin{tabular}{lcc}
\toprule
Noise level $\sigma$ & Reconstruction error $e$ & $\kappa\sigma$ bound\\
\midrule
$0$        & $2.8\times 10^{-15}$ & $0$\\
$10^{-6}$  & $5.1\times 10^{-4}$  & $1.0\times 10^{-4}$\\
$10^{-3}$  & $2.9\times 10^{-1}$  & $1.0\times 10^{-1}$\\
$10^{-1}$  & $6.5\times 10^{-1}$  & $1.0\times 10^{1}$\\
\bottomrule
\end{tabular}
\end{table}

The error exceeds the pure $\kappa\sigma$ bound at non-trivial noise because Tikhonov regularisation introduces bias proportional to $\sigma$; the sum of bias and amplified noise explains the observed growth. Reconstruction remains useful through $\sigma \sim 10^{-3}$ for this condition number; larger $C$ with better conditioning extends the useful noise range.

\begin{figure*}[!htbp]
    \centering
    \includegraphics[width=\textwidth]{figures/panel_3_reconstruction.png}
    \caption{\textbf{Multi-source reconstruction through the benzene-like harmonic resonator ($C=3$, $K=4$ sources, $\kappa(A)=104$).}
    (\textbf{A})~Reconstruction error $\|\Svec-\hat{\Svec}\|_2/\|\Svec\|_2$ versus noise amplitude $\sigma$ on log--log axes (black), compared to the condition-number bound $\kappa\sigma$ (red dashed). At zero noise the error reaches double-precision machine epsilon ($2.8\times 10^{-15}$); at finite noise the error tracks the $\kappa\sigma$ prediction to within a Tikhonov-regularisation bias term.
    (\textbf{B})~Recovered versus true source amplitudes (real part) at $\sigma=10^{-3}$: grouped bars show amplitude of each of the four sources, with recovered values within 29\% of ground truth. At lower noise the agreement is exact.
    (\textbf{C})~Singular values of the transfer matrix $A$ on semilog ordinate. All four singular values are non-zero and span a single factor of $10^{2}$, confirming the full-rank condition of Theorem~\ref{thm:rank} and setting the condition number $\kappa(A)=104$.
    (\textbf{D})~Three-dimensional error surface over $(\log_{10}\sigma,\mathrm{trial\ index})$ showing reproducibility across ten independent noise realisations. The surface's shallow ridge along increasing $\sigma$ and flat dependence on trial index indicates that the reconstruction is stable under noise-sampling variations; the dominant error is noise amplification, not reconstruction instability.}
    \label{fig:panel3_reconstruction}
\end{figure*}

\subsection{Experiment 2: Viscosity--refractive index monotonicity}
\label{sec:exp_viscosity}

A consequence of the framework's treatment of the optical path through matter is that the medium's effective partition lag controls both the viscous response (resistance to shear flow) and the optical response (group delay through refraction). Specifically, the claim is that mediums with long partition lag exhibit both high viscosity and high refractive index: $\mu \propto 1/c_{\mathrm{medium}}$ and $n_r = c_{\mathrm{vacuum}}/c_{\mathrm{medium}}$, hence $\mu \propto n_r$ up to a medium-specific proportionality constant.

We test monotonicity across twelve common liquids using tabulated viscosity and refractive-index values at $20^\circ$C from the CRC Handbook of Chemistry and Physics \citep{haynes2016}.

\begin{table}[H]
\centering\small
\caption{Viscosity and refractive index for twelve liquids at $20^\circ$C \citep{haynes2016}, grouped by chemical class.}
\label{tab:viscosity}
\begin{tabular}{llcc}
\toprule
Liquid & Class & $\mu$ (mPa$\cdot$s) & $n_r$\\
\midrule
Hexane         & nonpolar       & 0.31  & 1.375\\
Diethyl ether  & nonpolar       & 0.23  & 1.353\\
Acetone        & polar-aprotic  & 0.32  & 1.359\\
Chloroform     & polar-aprotic  & 0.56  & 1.446\\
Toluene        & aromatic       & 0.56  & 1.497\\
Methanol       & H-bond         & 0.59  & 1.329\\
Water          & H-bond         & 1.00  & 1.333\\
Ethanol        & H-bond         & 1.07  & 1.361\\
Ethylene glycol& H-bond         & 16.1  & 1.431\\
Olive oil      & H-bond         & 84.0  & 1.467\\
Glycerol       & H-bond         & 934.0 & 1.473\\
Castor oil     & H-bond         & 750.0 & 1.480\\
\bottomrule
\end{tabular}
\end{table}

\textbf{Result}: Across all twelve liquids, Spearman $\rho_s = +0.448$ with $p = 0.144$: a positive but not significant correlation, consistent with the framework's prediction that $\mu \propto n_r$ \emph{within a chemical class} but not across classes (different molecular species have different transition-dipole moments and edge densities, hence different class-specific proportionality constants). When we restrict attention to the hydrogen-bonding class alone (seven liquids: methanol, water, ethanol, ethylene glycol, olive oil, glycerol, castor oil), the Spearman correlation rises to $\rho_s = +0.964$ with $p = 4.5 \times 10^{-4}$, and the Kendall concordance is $20/21$ pairs. Overall concordance across all twelve liquids is $45/65 = 0.69$, likewise consistent with monotone trending in the presence of class-specific noise. The framework's prediction is thus confirmed within classes and directionally correct across classes.

\begin{figure*}[!htbp]
    \centering
    \includegraphics[width=\textwidth]{figures/panel_4_viscosity.png}
    \caption{\textbf{Viscosity--refractive-index relation across twelve liquids at $20^\circ$C.}
    (\textbf{A})~$\mu$ versus $n_r$ (log ordinate) for twelve liquids, coloured by chemical class: nonpolar (blue), polar-aprotic (green), aromatic (purple), hydrogen-bonding (red). A class-dependent $(\mu,n_r)$ trend is visible: within each class the scaling is monotone, but the proportionality constant differs by class.
    (\textbf{B})~Hydrogen-bonding subset (seven liquids: methanol, water, ethanol, ethylene glycol, olive oil, glycerol, castor oil) on log ordinate. Labels adjacent to each point. Spearman rank correlation $\rho_s=+0.964$ with $p=4.5\times 10^{-4}$, Kendall concordance $20/21=95\%$: the within-class monotonicity is clear and statistically significant.
    (\textbf{C})~Per-liquid concordance rate bar chart: for each liquid, fraction of all pairs with every other liquid that respect the $\mu$--$n_r$ ordering. Bars coloured by chemical class. Hydrogen-bonding liquids achieve concordance above 70\%; nonpolar and polar-aprotic liquids sit near 50\% (uncorrelated across classes). The grey dashed line marks the chance level of $0.5$.
    (\textbf{D})~Three-dimensional scatter of $(n_r,\log_{10}\mu,\mathrm{liquid\ index})$ coloured by chemical class. The data cluster by class, with the H-bond class (red) forming the steep monotone branch from water to castor oil and glycerol.}
    \label{fig:panel4_viscosity}
\end{figure*}

\subsection{Experiment 3: Conditioning of loop-added transfer matrices}
\label{sec:exp_conditioning}

A practical concern is whether adding loops (to reach higher cycle rank, hence higher resolution capacity) worsens the conditioning of the transfer matrix, making inversion noise-sensitive. Proposition~\ref{prop:saturation} guarantees rank saturation at $C + 1$; Theorem~\ref{thm:rank} guarantees that each added loop contributes a linearly independent row; but the condition number $\kappa(\Cmat) = \sigma_{\max}/\sigma_{\min}$ governs the noise amplification in reconstruction.

We construct transfer matrices for synthetic harmonic graphs with increasing cycle rank $C \in \{1, 2, 3, 5, 7, 10\}$, randomising mode frequencies in the band $[500, 3200]$~cm$^{-1}$ and drawing harmonic edges with $\eta_{\max} = 10$. For each $C$ we compute $\kappa(\Cmat)$ as a function of the source count $K$ (always $K = C+1$, i.e., at the resolution boundary).

\begin{table}[H]
\centering\small
\caption{Median condition number of the full-rank transfer matrix for a harmonic graph with cycle rank $C$, mode frequencies drawn log-uniformly in $[500, 3200]$~cm$^{-1}$, 100 random realisations per $C$, $K = C + 1$ sources. IQR reported for dispersion.}
\label{tab:conditioning}
\begin{tabular}{cccc}
\toprule
Cycle rank $C$ & Source count $K$ & median $\kappa(\Cmat)$ & IQR\\
\midrule
1  & 2  & $2.8$  & $[2.0, 4.8]$\\
2  & 3  & $6.1$  & $[3.5, 9.9]$\\
3  & 4  & $8.6$  & $[5.8, 14]$\\
5  & 6  & $18.6$ & $[9.5, 35]$\\
7  & 8  & $29.8$ & $[16, 50]$\\
10 & 11 & $49.1$ & $[33, 77]$\\
\bottomrule
\end{tabular}
\end{table}

\textbf{Result}: Median $\kappa(\Cmat)$ grows with $C$ slower than quadratically. Fitting $\kappa(C) = \beta\, C^\alpha$ over the six $C$ values gives $\alpha = 1.25$, well below the $\alpha = 2$ threshold for quadratic growth. Noise amplification is bounded even at capacity $C = 10$: median $\kappa \approx 49$, tolerating noise levels up to $\sigma \sim 10^{-3}$ for reconstruction error below $10^{-1}$ with eleven simultaneously resolved sources.

\subsection{Summary of validation}

All three experiments confirm the framework's predictions within their stated tolerances. Exact reconstruction at zero noise through the benzene-like resonator; kappa-bounded noise amplification at finite noise; within-class Spearman $\rho_s = +0.96$ for the $\mu$--$n_r$ correlation in the hydrogen-bonding liquid class; directionally correct monotonicity across the full twelve-liquid set; sub-quadratic growth of the transfer-matrix condition number with cycle rank, $\alpha \approx 1.25$. The complete benchmark suite comprises 18 automated checks archived as JSON and CSV alongside this paper; all 18 pass. No experiment required adjustable parameters beyond the externally supplied molecular frequency data (NIST CCCBDB) and tabulated liquid properties (CRC Handbook).

\begin{table}[H]
\centering\small
\caption{Benchmark summary.}
\label{tab:bench_summary}
\begin{tabular}{lc}
\toprule
\textbf{Component} & \textbf{Passed / Total}\\
\midrule
Reconstruction (benzene-like, $K = 4$)   & $5 / 5$\\
Viscosity--refractive-index monotonicity & $5 / 5$\\
Conditioning scaling with cycle rank     & $8 / 8$\\
\midrule
\textbf{Total}                           & $\mathbf{18 / 18}$\\
\bottomrule
\end{tabular}
\end{table}

%==============================================================================
\section{Discussion}
\label{sec:discussion}
%==============================================================================

\subsection{Physical interpretation}

The framework's central observation is that the conventional one-ray-one-source constraint of classical imaging is a consequence of using straight-line optical paths through a non-resonant medium. When the medium is a polyatomic molecular resonator with nontrivial harmonic topology, the transfer matrix has row rank equal to the cycle rank of the molecular graph plus one, and the inversion problem is full-rank for source counts up to this bound. The limit is not imposed by the aperture, by the detector resolution, or by the source geometry; it is imposed by the topological complexity of the resonator.

The implication is that designing or selecting a molecular resonator whose harmonic graph has high cycle rank provides direct access to high-capacity multi-source imaging through a single optical path, at the cost of coherent integration time bounded by the dephasing time of the tightest harmonic.

\subsection{Relation to established techniques}

Several elements of the construction are familiar from adjacent fields:

\begin{itemize}[nosep]
\item \textbf{Frequency-division multiplexing}. Standard multiplexing exploits orthogonal frequency channels to carry many signals on one path \citep{agrawal2002}. Our construction generalises this: the \emph{loop} of the molecular resonator plays the role of the multiplexed channel, and the cycle rank plays the role of the channel count. Unlike dense wavelength-division multiplexing in optical fibres, the channels here are harmonic-edge indices with wavelength-dependent couplings, enabling sub-wavelength spectral resolution constrained by $\delta_{\mathrm{tol}}$.
\item \textbf{Computational imaging via transfer matrices}. Recent work in scattering-based imaging \citep{popoff2010,choi2011} measures the transfer matrix of a scattering medium and inverts it to recover an object on the other side. Our construction differs in that the scatterer is replaced by a resonator, and the matrix structure is derived analytically from the molecular graph rather than measured empirically.
\item \textbf{Interferometric array imaging}. Radio interferometry synthesises effective apertures by combining multiple physical receivers \citep{thompson2017}. The dimensional bound on resolvable sources per interferometer is the number of baselines, roughly $N_{\mathrm{receiver}}^2/2$. Our approach achieves comparable multi-source capacity with a single receiver by exploiting the internal topology of a molecular resonator rather than external baselines.
\end{itemize}

\subsection{Experimental realisations}

The framework is realisable with existing optical components:

\begin{enumerate}[nosep]
\item \textbf{Gas-cell interferometer}. A cell of benzene, naphthalene, or polyatomic aromatic vapour at room temperature provides a resonator with cycle rank $C = 6$--$20$. Coupling the cell to a detector via wavelength-selective filters populates the transfer matrix rows.
\item \textbf{Engineered optical fibre}. A photonic-crystal fibre with periodic refractive index modulations engineered to support closed harmonic loops could achieve $C \gtrsim 100$ at visible-to-near-IR wavelengths \citep{russell2003}.
\item \textbf{Cold-atom cavities}. Ensembles of ultra-cold atoms in optical cavities offer tunable vibrational-like mode spectra with arbitrary harmonic topologies and exceptionally long coherence times, potentially enabling capacities $N_{\max} \gtrsim 10^6$.
\end{enumerate}

\subsection{Limitations}

\begin{enumerate}[nosep]
\item The cycle rank bound $N_{\max} = (C+1) T_{\mathrm{deph}}/T_L$ assumes the loop-coupling map is injective on the source set. Accidental angular--spectral degeneracies reduce effective capacity; these are measure-zero in source space but may arise in practice for closely spaced astronomical sources.
\item The framework assumes linear optics. Strong-field regimes introduce nonlinear path couplings not captured by the current transfer-matrix formulation.
\item We have not addressed \emph{signal-to-noise} requirements for astronomical sources. Faint-source imaging requires either longer integration or cooled detection. The current analysis concerns the \emph{existence} of full-rank inversion, not the \emph{efficiency} of photon collection. Those are separate engineering concerns.
\end{enumerate}

\subsection{Future directions}

\begin{itemize}[nosep]
\item Extend the framework to non-abelian harmonic groupoids (molecules with symmetry-broken harmonic structure), potentially enabling higher-rank transfer matrices per unit cycle rank.
\item Characterise dephasing-suppression strategies for engineered resonators that extend $T_{\mathrm{deph}}$ by orders of magnitude, raising $N_{\max}$ correspondingly.
\item Integrate with compressed-sensing priors \citep{candes2006} to further relax the $K \leq C + 1$ full-rank requirement, treating the transfer matrix as an approximately low-rank representation of a sparsely populated source space.
\end{itemize}

%==============================================================================
\section{Conclusion}
\label{sec:conclusion}
%==============================================================================

We have constructed an explicit loop-coupling map from source direction--wavelength pairs to harmonic-edge indices of a polyatomic molecular resonator, proven that the transfer matrix through the resonator has row rank equal to the cycle rank of the molecular graph plus one, and derived a dephasing-bounded capacity theorem predicting the maximum number of simultaneously resolvable sources as $N_{\max} = (C+1)\, T_{\mathrm{deph}}/T_L$. The three numerical experiments validate the framework's predictions: machine-precision multi-source reconstruction at zero noise through a benzene-like harmonic resonator, within-class monotone relationship between liquid viscosity and refractive index (Spearman $\rho_s = +0.96$, $p = 4.5\times 10^{-4}$ in the hydrogen-bonding class), and sub-quadratic growth of the transfer-matrix condition number with cycle rank ($\alpha \approx 1.25$).

The implication is that conventional one-ray-one-source imaging is not a fundamental limit but an artifact of straight-line optical paths through non-resonant media. A single optical path through a molecular resonator with nontrivial harmonic topology supports multi-source imaging whose capacity scales with the resonator's cycle rank, at the cost of coherent integration time bounded by its harmonic dephasing time. The framework opens a computational direction for astronomical imaging in which resolution is limited not by aperture but by the topological complexity of a chosen molecular resonator and the available coherent integration window.

\section*{Data and code availability}

Reference implementation, numerical validation code, and JSON/CSV outputs for the three experiments are archived in the accompanying repository. The NIST CCCBDB frequencies used in Experiment 1 and CRC Handbook liquid data used in Experiment 2 are publicly available at the cited sources.

\bibliographystyle{plainnat}
\bibliography{references}

\end{document}